\subsection{Трение на горке}
%%%%%%%%%%%%%%%%%%%%%%%%%%%%%%%%%%%%%%%%%%%%%%%%%%%%
%%%%%%%%%%%%%%%%%%%%%%%%%%%%%%%%%%%%%%%%%%%%%%%%%%%%
%%%%%%%%%%%%%%%%%%%%%%%%%%%%%%%%%%%%%%%%%%%%%%%%%%%%	
 \z[{\begin{wrapfigure}[5]{r}{0.3\textwidth}%
		\vspace{-10pt}%
		\begin{flushright}%
 			\includegraphics[width=0.3\textwidth]{5.3.1}%
		\end{flushright}%
	\end{wrapfigure}%
	}]%
	{На гладкой горизонтальной поверхности лежит клин с углом $\alpha = 15\degree$ при основании, упираясь торцом в неподвижную вертикальную стенку. По верхней грани клина соскальзывает без трения тело массы $m = 0.2\,\text{кг}$. Найдите силу нормального давления $N$ клина на стенку.}
 %
{$N=mg\sin \alpha = 5.65\,\text{Н}$}
%%%%%%%%%%%%%%%%%%%%%%%%%%%%%%%%%%%%%%%%%%%%%%%%%%%%	
 \z[{\begin{wrapfigure}[5]{r}{0.3\textwidth}%
		\vspace{-12pt}%
		\begin{flushright}%
 			\includegraphics[width=0.3\textwidth]{5.3.2}%
		\end{flushright}%
	\end{wrapfigure}%
	}]%
	{ Человек массой $m$, упираясь ногами в ящик массой $M$, подтягивает его с помощью каната, перекинутого через блок, по наклонной плоскости с углом наклона $\alpha$. С какой минимальной силой надо тянуть канат человеку, чтобы подтянуть ящик к блоку? Коэффициент трения между ящиком и наклонной плоскостью равен $\mu$.}
 %
{$T_{\min}=\frac{m+M}{2}g(\mu \cos \alpha + \sin \alpha)$}
%%%%%%%%%%%%%%%%%%%%%%%%%%%%%%%%%%%%%%%%%%%%%%%%%%%% 
 \z[{\begin{wrapfigure}[5]{r}{0.3\textwidth}%
		\vspace{-12pt}%
		\begin{flushright}%
 			\includegraphics[width=0.3\textwidth]{5.3.3}%
		\end{flushright}%
	\end{wrapfigure}%
	}]%
	{Гладкий сверху, но шероховайтый снизу клин (угол $\alpha = \pi/6$) массой $M$ находится на шероховатой горизонтальной поверхности стола. Через блок, укреплённый на вершине клина, перекинута лёгкая нерастяжимая нить, связывающая грузы, массы которых $m_1$ и $m_2$. Грузы удерживают, затем отпускают. После этого грузы движутся, а клин покоится. Считайте $M = 2m$, $m_1  = m$, $m_2  = 2m$. Найдите: 
 	\lpu ускорение грузов;\lpuend
 	\lpu силу $T$ натяжения нити; \lpuend
 	\lpu силу $N$ нормальной реакции, действующей на клин со стороны стола.\lpuend
	}
%
{\puan $a=\frac{2-\sin \alpha}{3}g;$~\puan $T=mg;$~\puan$N=\frac{17}{4}mg$}
%%%%%%%%%%%%%%%%%%%%%%%%%%%%%%%%%%%%%%%%%%%%%%%%%%%% 
 \z[{\begin{wrapfigure}[5]{r}{0.3\textwidth}%
		\vspace{-12pt}%
		\begin{flushright}%
 			\includegraphics[width=0.3\textwidth]{5.3.4}%
		\end{flushright}%
	\end{wrapfigure}%
	}]%
	{Бруски с массами $m$ и $2m$ связаны лёгкой нитью, перекинутой через блок, и находятся на наклонной и горизонтальной поверхностях призмы. Угол наклона к горизонту одной из поверхностей призмы равен $\alpha\,(\sin \alpha = 3/5)$. Коэффициент трения скольжения бруска о горизонтальную поверхность $ \mu = 1/6$, а о наклонную поверхность -- $2\mu$. При перемещении призмыс некоторым минимальным горизонтальным ускорением $a$ брусок с массой $2m$ начинает скользить по призме влевопри натянутой нити. Найти ускорение $a$.}
 %
{$a=\frac{6}{13}g$}
%%%%%%%%%%%%%%%%%%%%%%%%%%%%%%%%%%%%%%%%%%%%%%%%%%%%
 \z[{\begin{wrapfigure}[7]{r}{0.3\textwidth}%
		\vspace{-15pt}%
		\begin{flushright}%
 			\includegraphics[width=0.3\textwidth]{5.3.5}%
		\end{flushright}%
	\end{wrapfigure}%
	}]%
	{На наклонной плоскости, образующей с горизонтальной поверхностью угол $\vartheta\ = 30\degree$, закреплен желоб, как показано на рисунке. С плоскостью связана координатная система $XOY$, начало которой совмещено с нижней точкой желоба. По желобу из состояния покоя начинает соскальзывать маленькая гирька. Найдите скорость $v$ гирьки в нижней точке желоба, если начальные координаты гирьки $x_0  = 0.5\,\text{м}$, $y_0=1\,\text{м}$, а коэффициент трения между гирькой и желобом $\mu = 0.3$.}
 %
{$v=2\,\text{м}/\text{с}$}